\documentclass[11pt,a4paper]{article}

% --- Packages ---
\usepackage[margin=.75in]{geometry}
\usepackage[utf8]{inputenc}
\usepackage[T1]{fontenc}
\usepackage{lmodern} % Fix for font not found error
\usepackage{titlesec}
\usepackage{enumitem}
\usepackage{hyperref}
\usepackage{xcolor}
\usepackage{fancyhdr}
\usepackage{comment}
\usepackage{cleveref}
\usepackage[bottom]{footmisc}
\usepackage{graphicx}
\usepackage{longtable}
\usepackage{booktabs}
\usepackage{array}
\usepackage{calc}
\usepackage{etoolbox}

% Standard Pandoc setup
\providecommand{\tightlist}{%
  \setlength{\itemsep}{0pt}\setlength{\parskip}{0pt}}

\crefformat{section}{\S#2#1#3}
\crefformat{subsection}{\S#2#1#3}
\crefformat{subsubsection}{\S#2#1#3}
\crefformat{figure}{#2Figure~#1#3}
\crefformat{footnote}{#2\footnotemark[#1]#3}

\linespread{0.9} % Riduce lo spazio tra le linee

% --- Colors ---
\definecolor{primary}{RGB}{0, 0, 0}
\definecolor{links}{HTML}{0000EE}
\definecolor{blue}{HTML}{26428B}

\hypersetup{
    colorlinks=true,
    linkcolor=links,
    urlcolor=links,
    citecolor=links,
}

% --- Section Formatting ---
\titleformat{\section}
  {\normalfont\Large\bfseries\color{blue}}{\thesection}{1em}{}
\titleformat{\subsection}
{\normalfont\large\bfseries\color{blue}}{\thesubsection}{1em}{}
\titleformat{\subsubsection}
{\normalfont\bfseries\color{blue}}{\thesubsubsection}{1em}{}

% --- Header and Footer ---
\pagestyle{fancy}
\fancyhead[L]{\small Giuseppe Capasso}
\fancyhead[R]{\small intro-to-linux}
\fancyfoot[C]{\thepage}
\renewcommand{\headrulewidth}{0.4pt}

% --- Custom Title Command ---
\newcommand{\proposaltitle}[1]{
    \begin{center}
        \vspace*{0.1cm}
        {\LARGE \bfseries \color{blue} #1} \\[1.5ex]
        {\large \textbf{Giuseppe Capasso}} \\[1ex]
                \small{
            \vspace{0.1cm}
            \textbf{Materia:} intro-to-linux\\
        }
            \end{center}
    \vspace{0.3cm}
    \hrule height 0.5pt
    \vspace{0.5cm}
}

\begin{document}

\proposaltitle{Laboratorio Pratico: Da Zero a SysAdmin}

Laboratorio Pratico: Da Zero a SysAdmin (Ubuntu Server) Obiettivo:
Installare un server Linux ``nudo'', prenderne il controllo da remoto
tramite Windows, metterlo in sicurezza con un firewall e creare un
servizio di sistema immortale. Strumenti: VirtualBox, Ubuntu Server
(ISO), Windows Terminal / PowerShell.

FASE 1: Creazione VM e ``Port Forwarding'' I veri server sono chiusi nei
datacenter. Per simulare questa condizione, la nostra VM sarà isolata in
modalità NAT, e creeremo dei ``tunnel'' per potervi accedere dal nostro
PC Windows fisico. 1. Crea la Macchina Virtuale in VirtualBox: • OS:
Ubuntu Server (64-bit). • RAM: 2048 MB (2 GB). • Disco: 15 GB. •
Attenzione durante l'installazione: Quando l'installer te lo chiede,
spunta la casella ``Install OpenSSH server''. Crea un utente chiamato
studente con password studente. 2. Configura il Port Forwarding (A VM
spenta o accesa): • Vai su Impostazioni della VM -\textgreater{} Rete
-\textgreater{} Avanzate -\textgreater{} Inoltro delle porte. • Crea due
regole cliccando sull'icona {[}+{]}: • Regola 1 (SSH): Nome SSH,
Protocollo TCP, IP Host 127.0.0.1, Porta Host 2222, Porta Guest 22. •
Regola 2 (Web): Nome Web, Protocollo TCP, IP Host 127.0.0.1, Porta Host
8080, Porta Guest 8000. • Clicca OK e avvia la VM. Fai il login una
volta sola dalla finestra nera di VirtualBox per assicurarti che sia
viva, poi riducila a icona. Non la useremo più.

FASE 2: Presa di controllo remota (Windows Terminal) Ora lavoreremo come
veri amministratori di sistema, dal terminale del nostro PC fisico. 1.
Apri PowerShell (o Windows Terminal) sul tuo PC.

2. Collegati al server usando il tunnel che abbiamo creato sulla porta
2222: PowerShell ssh -p 2222 studente@127.0.0.1

(Scrivi yes per accettare il fingerprint e inserisci la password
studente).

2.1 Configurazione Autenticazione a Chiavi (SSH Keys) Basta password.
Usiamo la crittografia. 1. Apri una SECONDA scheda di PowerShell (non
dentro Linux, ma sul tuo Windows locale). 2. Genera le chiavi (premi
sempre Invio per lasciare i valori di default): PowerShell ssh-keygen -t
ed25519

\begin{enumerate}
\def\labelenumi{\arabic{enumi}.}
\setcounter{enumi}{2}
\item
  Copia la chiave pubblica sul server Linux usando questo comando magico
  da PowerShell: PowerShell cat \$env:USERPROFILE.ssh\id\_ed25519.pub
  \textbar{} ssh -p 2222 studente@127.0.0.1 ``mkdir -p
  \textasciitilde/.ssh \&\& cat \textgreater\textgreater{}
  \textasciitilde/.ssh/authorized\_keys''
\item
  Test: Torna alla prima scheda (quella collegata a Linux), scrivi exit
  per uscire, e poi ricollegati con ssh -p 2222 studente@127.0.0.1. Se
  entri istantaneamente senza che ti venga chiesta la password, ha
  funzionato!
\end{enumerate}

2.2 Blindare SSH Ora che hai la chiave, chiudi la porta a chi tenta di
usare le password. 1. Nel terminale Linux, modifica il file di
configurazione SSH: Bash sudo nano /etc/ssh/sshd\_config

\begin{enumerate}
\def\labelenumi{\arabic{enumi}.}
\setcounter{enumi}{1}
\item
  Trova la riga \#PasswordAuthentication yes, togli il commento (\#) e
  cambiala in: Plaintext PasswordAuthentication no
\item
  Salva (CTRL+O, CTRL+X) e riavvia il servizio: Bash sudo systemctl
  restart ssh
\end{enumerate}

FASE 3: Sicurezza (UFW Firewall) Il server è raggiungibile, ora dobbiamo
difenderlo. Attenzione all'ordine dei comandi o ti chiuderai fuori! 1.
Blocca tutto il traffico in entrata di default: Bash sudo ufw default
deny incoming

\begin{enumerate}
\def\labelenumi{\arabic{enumi}.}
\setcounter{enumi}{1}
\item
  Consenti il traffico in uscita: Bash sudo ufw default allow outgoing
\item
  VITALIZIO: Apri la porta SSH prima di attivare il firewall! Bash sudo
  ufw allow 22/tcp
\item
  Apri la porta per il nostro futuro server web Python (porta 8000):
  Bash sudo ufw allow 8000/tcp
\item
  Attiva il firewall: Bash sudo ufw enable
\end{enumerate}

(Rispondi y alla domanda se vuoi procedere). 6. Verifica che le regole
siano corrette: Bash sudo ufw status

FASE 4: Systemd - Il Servizio Immortale Costruiamo un piccolo server web
e diciamo a Linux di tenerlo sempre acceso in background. 1. Prepara i
file del sito: Crea una directory e un file HTML di base. Bash sudo
mkdir -p /var/www/pythonweb echo ``

Vittoria! Il servizio Systemd funziona!

'' \textbar{} sudo tee /var/www/pythonweb/index.html

2. Crea il file del Servizio (Unit File): Questo file spiega a Systemd
come far partire la nostra applicazione. Bash sudo nano
/etc/systemd/system/pythonweb.service

\begin{enumerate}
\def\labelenumi{\arabic{enumi}.}
\setcounter{enumi}{2}
\tightlist
\item
  Incolla questa configurazione esatta: Ini, TOML {[}Unit{]}
  Description=Web Server Python di Test After=network.target
\end{enumerate}

{[}Service{]} Type=simple User=studente
WorkingDirectory=/var/www/pythonweb ExecStart=/usr/bin/python3 -m
http.server 8000 Restart=always RestartSec=3

{[}Install{]} WantedBy=multi-user.target

(Salva con CTRL+O ed esci con CTRL+X). 4. Inizializza e Avvia il
Servizio: Invia i comandi a Systemd per caricare il file, avviare il
server e abilitarlo al boot automatico: Bash sudo systemctl
daemon-reload sudo systemctl start pythonweb sudo systemctl enable
pythonweb sudo systemctl status pythonweb

(Dovresti vedere un pallino verde e la scritta active (running). Premi q
per uscire dalla schermata di status).

FASE 5: Il Test Finale (Prova del Nove) È il momento di raccogliere i
frutti del lavoro. 1. Test del Web Server: Apri Google Chrome o Edge sul
tuo PC Windows. Nella barra degli indirizzi, digita:
http://127.0.0.1:8080 (Se vedi la scritta ``Vittoria!'', il port
forwarding e il servizio Systemd stanno dialogando perfettamente). 2.
Test di Resilienza (Chaos Monkey): Torna nel terminale Linux. Uccidiamo
brutalmente il processo Python per simulare un crash del server: Bash
sudo killall python3

Ora ricarica velocemente la pagina web su Windows. Funziona ancora?
Controlla lo stato del servizio: Bash sudo systemctl status pythonweb

Noterai che Systemd, obbedendo alla direttiva Restart=always, ha
rianimato il processo in meno di 3 secondi!



\end{document}
